% GENERATED FILE. Edit canonical lessons, then run npm run generate:books.
% canonical-source-hash: fnv1a64:5631139bd165861a
% canonical-lessons: LA-C202-suadeo, LA-C202-tremo, LA-C202-excipio, LA-C202-ago, LA-C202-procedo

\chapter{To advise, To tremble, To welcome, To do, To go forward}
\label{ch:la-to-advise-to-tremble-to-welcome-to-do-to-go-forward}
\hlchaptermodality{202}

\begin{chapteropening}
\textbf{By the end of this chapter:} \emph{I can say to advise, to tremble, to welcome, to do and to go forward.}

Complete the last lesson of chapter 202: I can say to advise, to tremble, to welcome, to do and to go forward.
\end{chapteropening}

\section[suādeō, suādēre]{suādeō, suādēre — to advise, to urge}

\label{lesson:LA-C202-suadeo}



\begin{quote}
Before the new one: say the Latin for to celebrate, then the Latin for to pursue.
\end{quote}

\subsection*{You'll want to know: suādeō, suādēre}
\textbf{suādeō, suādēre} — \textquotedblleft{}to advise, to urge\textquotedblright{}.

\begin{cousinweb}
\textbf{suādeō} is inherited, usually traced to Proto-Indo-European \emph{*sweh2d-}, \textquotedblleft{}sweet\textquotedblright{} (EtymDB 2.1): to advise was to make a thing sweet to someone. \emph{suāvis}, \textquotedblleft{}sweet\textquotedblright{}, and English \emph{sweet} are cousins (Lewis \& Short).
\end{cousinweb}

\begin{grammarlens}[title={What you've built}]
One more word for this chapter.
\end{grammarlens}

\subsection*{Your turn}
\begin{itemize}
\raggedright
  \item \emph{Say it:} \emph{suādeō, suādēre}
  \item \emph{Say it:} \emph{suādeō, suādēre}, once more
  \item \emph{Say it:} say \emph{suādeō, suādēre}
  \item \emph{Recall:} say the Latin for to celebrate, then the Latin for to pursue, then say \emph{suādeō, suādēre} again
\end{itemize}

\begin{tcolorbox}[breakable,colback=teal!4,colframe=teal!35!black,title={Before you move on}]
What does suādeō, suādēre mean? (\textquotedblleft{}To advise, to urge\textquotedblright{}.) Say it once more.
\end{tcolorbox}
\section[tremō, tremere]{tremō, tremere — to tremble}

\label{lesson:LA-C202-tremo}



\begin{quote}
Before the new one: say the Latin for to pursue, then the Latin for to advise.
\end{quote}

\subsection*{You'll want to know: tremō, tremere}
\textbf{tremō, tremere} — \textquotedblleft{}to tremble\textquotedblright{}.

\begin{cousinweb}
\textbf{tremō} is inherited, usually traced to Proto-Indo-European \emph{*trem-}, \textquotedblleft{}to tremble\textquotedblright{} (EtymDB 2.1); Greek \emph{tremō}, \textquotedblleft{}to tremble\textquotedblright{}, is a cousin (Lewis \& Short).
\end{cousinweb}

\begin{grammarlens}[title={What you've built}]
One more word for this chapter.
\end{grammarlens}

\subsection*{Your turn}
\begin{itemize}
\raggedright
  \item \emph{Say it:} \emph{tremō, tremere}
  \item \emph{Say it:} \emph{tremō, tremere}, once more
  \item \emph{Say it:} say \emph{tremō, tremere}
  \item \emph{Recall:} say the Latin for to pursue, then the Latin for to advise, then say \emph{tremō, tremere} again
\end{itemize}

\begin{tcolorbox}[breakable,colback=teal!4,colframe=teal!35!black,title={Before you move on}]
What does tremō, tremere mean? (\textquotedblleft{}To tremble\textquotedblright{}.) Say it once more.
\end{tcolorbox}
\section[excipiō, excipere]{excipiō, excipere — to welcome, to receive}

\label{lesson:LA-C202-excipio}



\begin{quote}
Before the new one: say the Latin for to advise, then the Latin for to tremble.
\end{quote}

\subsection*{You'll want to know: excipiō, excipere}
\textbf{excipiō, excipere} — \textquotedblleft{}to welcome, to receive\textquotedblright{}.

\begin{cousinweb}
\textbf{excipiō} is \emph{ex}, \textquotedblleft{}out\textquotedblright{}, with \emph{capiō}, \textquotedblleft{}to take\textquotedblright{} (EtymDB 2.1), its a weakened to i: to take out, to receive. \emph{capiō} is usually traced to Proto-Indo-European \emph{*keh2p-}, \textquotedblleft{}to grab\textquotedblright{} (EtymDB 2.1).
\end{cousinweb}

\begin{grammarlens}[title={What you've built}]
One more word for this chapter.
\end{grammarlens}

\subsection*{Your turn}
\begin{itemize}
\raggedright
  \item \emph{Say it:} \emph{excipiō, excipere}
  \item \emph{Say it:} \emph{excipiō, excipere}, once more
  \item \emph{Say it:} say \emph{excipiō, excipere}
  \item \emph{Recall:} say the Latin for to advise, then the Latin for to tremble, then say \emph{excipiō, excipere} again
\end{itemize}

\begin{tcolorbox}[breakable,colback=teal!4,colframe=teal!35!black,title={Before you move on}]
What does excipiō, excipere mean? (\textquotedblleft{}To welcome, to receive\textquotedblright{}.) Say it once more.
\end{tcolorbox}
\section[agō, agere]{agō, agere — to do, to drive}

\label{lesson:LA-C202-ago}



\begin{quote}
Before the new one: say the Latin for to tremble, then the Latin for to welcome.
\end{quote}

\subsection*{You'll want to know: agō, agere}
\textbf{agō, agere} — \textquotedblleft{}to do, to drive\textquotedblright{}.

\begin{cousinweb}
\textbf{agō} is inherited, usually traced to Proto-Indo-European \emph{*h2eg-}, \textquotedblleft{}to drive, to lead\textquotedblright{} (EtymDB 2.1); Greek \emph{agō}, \textquotedblleft{}to lead\textquotedblright{}, is a cousin.
\end{cousinweb}

\begin{grammarlens}[title={What you've built}]
One more word for this chapter.
\end{grammarlens}

\subsection*{Your turn}
\begin{itemize}
\raggedright
  \item \emph{Say it:} \emph{agō, agere}
  \item \emph{Say it:} \emph{agō, agere}, once more
  \item \emph{Say it:} say \emph{agō, agere}
  \item \emph{Recall:} say the Latin for to tremble, then the Latin for to welcome, then say \emph{agō, agere} again
\end{itemize}

\begin{tcolorbox}[breakable,colback=teal!4,colframe=teal!35!black,title={Before you move on}]
What does agō, agere mean? (\textquotedblleft{}To do, to drive\textquotedblright{}.) Say it once more.
\end{tcolorbox}
\section[prōcēdō, prōcēdere]{prōcēdō, prōcēdere — to go forward, to advance}

\label{lesson:LA-C202-procedo}



\begin{quote}
Before the new one: say the Latin for to welcome, then the Latin for to do.
\end{quote}

\subsection*{You'll want to know: prōcēdō, prōcēdere}
\textbf{prōcēdō, prōcēdere} — \textquotedblleft{}to go forward, to advance\textquotedblright{}.

\begin{cousinweb}
\textbf{prōcēdō} is \emph{prō}, \textquotedblleft{}forward\textquotedblright{}, with \emph{cēdō}, \textquotedblleft{}to go\textquotedblright{} (EtymDB 2.1). English \emph{proceed} and \emph{process} come from it.
\end{cousinweb}

\begin{grammarlens}[title={What you've built}]
One more word for this chapter.
\end{grammarlens}

\subsection*{Your turn}
\begin{itemize}
\raggedright
  \item \emph{Say it:} \emph{prōcēdō, prōcēdere}
  \item \emph{Say it:} \emph{prōcēdō, prōcēdere}, once more
  \item \emph{Say it:} say \emph{prōcēdō, prōcēdere}
  \item \emph{Recall:} say the Latin for to welcome, then the Latin for to do, then say \emph{prōcēdō, prōcēdere} again
\end{itemize}

\begin{tcolorbox}[breakable,colback=teal!4,colframe=teal!35!black,title={Before you move on}]
What does prōcēdō, prōcēdere mean? (\textquotedblleft{}To go forward, to advance\textquotedblright{}.) Say it once more.
\end{tcolorbox}
